\section{Standard Rules}\label{sec:standard_rules}

This appendix collects the standard meta-functions and rules of the
Filtered Stepper Calculus (substitution and composition) deferred from \autoref{sec:filter}. The syntax of
evaluation contexts and the value judgment are given in
\autoref{fig:filter-syntax} and the accompanying rules in the main
body.

\fbox{\([v / x] e = e'\)} \(e'\) can be obtained by substitution of \(v\) for
\(x\) in expression \(e\), renaming bound variables as needed so that
no free variable of \(v\) is captured.
\[
  \begin{aligned}
    [v / x] \Nat{n} &= \Nat{n} \\
    [v / x] x &= v \\
    [v / x] y &= y && \text{if } x \neq y \\
    [v / x] (e_1(e_2)) &= ([v / x] e_1)([v / x] e_2) \\
    [v / x] (e_1 + e_2) &= ([v / x] e_1) + ([v / x] e_2) \\
    [v / x] \Lam{y}{e} &= \Lam{y}{[v / x] e} && \text{if } x \neq y \\
    [v / x] \Lam{y}{e} &= \Lam{y}{e} && \text{if } x = y \\
    [v / x] \Filter{(p, a, g)}{e} &= \Filter{([v / x]p, a, g)}{[v / x] e} \\
    [v / x] \Residue{a}{l}{e} &= \Residue{a}{l}{[v / x] e} \\
    [v / x] \Fix{x}{e} &= \Fix{x}{e} \\
    [v / x] \Fix{y}{e} &= \Fix{y}{[v / x]e} && \text{if } x \neq y
  \end{aligned}
\]

\fbox{\([v / x] p = p'\)} \(p'\) can be obtained by the same substitution on
pattern \(p\), again renaming bound variables to avoid capture.
\[
  \begin{aligned}
    [v / x] \$e &= \$e \\
    [v / x] \$v &= \$v \\
    [v / x] \Nat{n} &= \Nat{n} \\
    [v / x] x &= \Patternize(v) \\
    [v / x] y &= y && \text{if } x \neq y \\
    [v / x] (p_1(p_2)) &= ([v / x] p_1)([v / x] p_2) \\
    [v / x] (p_1 + p_2) &= ([v / x] p_1) + ([v / x] p_2) \\
    [v / x] \Lam{y}{p} &= \Lam{y}{[v / x] p} && \text{if } x \neq y \\
    [v / x] \Lam{y}{p} &= \Lam{y}{p} && \text{if } x = y \\
    [v / x] \Fix{x}{p} &= \Fix{x}{p} \\
    [v / x] \Fix{y}{p} &= \Fix{y}{[v / x]p} && \text{if } x \neq y
  \end{aligned}
\]

\fbox{\(\Patternize(e) = p\)} \(p\) is the pattern, without wildcards,
that describes the stripped expression \(e\).
\[
  \begin{aligned}
    \Patternize(x) &= x \\
    \Patternize(\Nat{n}) &= \Nat{n} \\
    \Patternize(e_1(e_2)) &= \Patternize(e_1)(\Patternize(e_2)) \\
    \Patternize(e_1 + e_2) &= \Patternize(e_1) + \Patternize(e_2) \\
    \Patternize(\Lam{x}{e}) &= \Lam{x}{\Patternize(e)} \\
    \Patternize(\Fix{x}{e}) &= \Fix{x}{\Patternize(e)} \\
    \Patternize(\Filter{f}{e}) &= \Patternize(e) \\
    \Patternize(\Residue{a}{l}{e}) &= \Patternize(e)
  \end{aligned}
\]
The mechanization proves
\(\Matches{\Patternize(e_1)}{e_2}\) if and only if
\(\Strip{e_1} \equiv_\alpha \Strip{e_2}\).

Residue-only stripping \(\StripResidues{e}\) has the same clauses as
\(\Strip{e}\), except that it keeps filters:
\(\StripResidues{\Filter{f}{e}} = \Filter{f}{\StripResidues{e}}\).

\begin{samepage}
\jbox{\(\Compose{e}{\mathcal{E}}{e'}\)} Expression \(e\) can be obtained by plugging \(e'\) into the mark of \(\mathcal{E}\).
We also write \(\Plug{\mathcal{E}}{e'}\) for this \(e\), and
\(\Plug{\mathcal{E}}{\mathcal{E}_0}\) for the context obtained by
plugging \(\mathcal{E}_0\) into the mark of \(\mathcal{E}\).
\begin{mathpar}
  \inferrule[C-Top]{
  }{
    \Compose{e}{\circ}{e}
  } \\
  \inferrule[C-Ap-L]{
    \Compose{e_1}{\mathcal{E}_1}{e_1'}
  }{
    \Compose{e_1(e_2)}{\mathcal{E}_1(e_2)}{e_1'}
  } \qquad
  \inferrule[C-Ap-R]{
    \Compose{e_2}{\mathcal{E}_2}{e_2'}
  }{
    \Compose{e_1(e_2)}{e_1(\mathcal{E}_2)}{e_2'}
  } \\
  \inferrule[C-Add-L]{
    \Compose{e_1}{\mathcal{E}_1}{e_1'}
  }{
    \Compose{e_1 + e_2}{\mathcal{E}_1 + e_2}{e_1'}
  } \qquad
  \inferrule[C-Add-R]{
    \Compose{e_2}{\mathcal{E}_2}{e_2'}
  }{
    \Compose{e_1 + e_2}{e_1 + \mathcal{E}_2}{e_2'}
  } \\
  \inferrule[C-Filter]{
    \Compose{e}{\mathcal{E}}{e'}
  }{
    \Compose{\Filter{f}{e}}{\Filter{f}{\mathcal{E}}}{e'}
  } \qquad
  \inferrule[C-Residue]{
    \Compose{e}{\mathcal{E}}{e'}
  }{
    \Compose{\Residue{a}{l}{e}}{\Residue{a}{l}{\mathcal{E}}}{e'}
  }
\end{mathpar}
\end{samepage}

%%% Local Variables:
%%% mode: LaTeX
%%% TeX-master: "main"
%%% End:
